%% EB-JEPA Full Research Paper
%% Timothéo Courtois, Théo Vasnier, Yassine Hammache — June 2026
\documentclass[11pt,a4paper]{article}
\usepackage[utf8]{inputenc}
\usepackage[T1]{fontenc}
\usepackage{amsmath,amssymb,bm,mathtools}
\usepackage{booktabs,multirow,array,tabularx,longtable}
\usepackage{graphicx}
\usepackage[colorlinks=true,citecolor=blue,linkcolor=blue,urlcolor=blue]{hyperref}
\usepackage{geometry}
\usepackage[numbers,sort&compress]{natbib}
\usepackage{xcolor}
\usepackage{caption}
\usepackage{subcaption}
\usepackage{enumitem}
\usepackage{setspace}
\usepackage{parskip}
\usepackage{colortbl}
\geometry{top=1in,bottom=1in,left=1.1in,right=1.1in}

\newcolumntype{C}[1]{>{\centering\arraybackslash}p{#1}}
\newcolumntype{L}[1]{>{\raggedright\arraybackslash}p{#1}}

\definecolor{ours}{HTML}{d4edda}
\definecolor{sota}{HTML}{ffeeba}
\definecolor{hilight}{HTML}{cce5ff}

\title{\LARGE\bfseries EB-JEPA: Energy-Based Joint Embedding Predictive Architecture\\
for Self-Supervised EEG Representation Learning\\[4pt]
{\large A Technical Research Report}}

\author{
  Yani Hammache$^{\dagger}$ \quad
  Th\'eo Basseras$^{\dagger}$ \quad
  Titouan Vasnier$^{\dagger}$ \quad
  Timothy Courtois$^{\dagger}$\\[4pt]
  \small$^{\dagger}$Hackathon Team --- \textit{Hack the World}, Vivatech 2026
}
\date{June 2026}

\begin{document}
\maketitle

%% ============================================================
\begin{abstract}
%% ============================================================
We present EB-JEPA, an Energy-Based Joint Embedding Predictive Architecture
for self-supervised learning of EEG representations applied to clinical abnormality
detection.  Starting from a single EEG window, two independently corrupted views
are produced through different augmentation pathways and mapped to the same latent
space by a shared encoder trained without labels.  Four encoder families are evaluated
on the TUH Abnormal EEG Corpus (TUAB v3.0.1~\cite{tuab}, 2\,717 training and 276 evaluation
recordings): a temporal Conv1D network (0.4M parameters), a patch Transformer inspired
by LaBraM~\cite{labram} (1.2M), an FFT-tokenised Transformer inspired by BIOT~\cite{biot}
(1.2M), and a hierarchical channel-pool Transformer inspired by EEGPT~\cite{eegpt} (0.8M).
We further introduce NaiveJEPA, a temporal-split predictive variant requiring no
anti-collapse regulariser, evaluated on both TUAB and the 6-class TUH EEG Events corpus
(TUEV v2.0.0) to assess cross-task transfer.

Our best frozen encoder (Conv1D + corruption augmentation + SIGReg) reaches
\textbf{0.775 BACC per-window} and \textbf{0.825 BACC per-recording} on TUAB,
matching published ContraWR and approaching LaBraM-Base (0.814).
The central contribution is \emph{systematic diagnosis}: we document three failure
modes --- an architectural conflict between EEGPT's channel-pooling and channel-drop
augmentation, a per-window vs.\ per-recording evaluation gap (${\approx}5$\,pp) that
initially caused an inflated performance claim, and a spectral auxiliary loss whose
apparent benefit (0.836 per-recording) was entirely due to a confounded evaluation.
We further report a novel finding: VICReg with its \emph{default} invariance coefficient
$\lambda{=}1$ outperforms the paper-correct $\lambda{=}25$ by $3\,\text{pp}$ (${\approx}6\sigma$,
3 seeds each) on this EEG task, suggesting that anti-collapse regularisation (variance +
covariance) should dominate invariance for short-window clinical signals.
\end{abstract}

\tableofcontents
\newpage

%% ============================================================
\section*{Introduction}
\addcontentsline{toc}{section}{Introduction}
\label{sec:intro}
%% ============================================================

\paragraph{Motivation.}
Clinical EEG abnormality detection is a high-stakes task: a board-certified neurologist
must review each recording to label it as normal or abnormal, a process that is
time-consuming and subject to inter-rater variability.  Self-supervised learning (SSL)
offers the prospect of learning useful representations from the vast majority of
unlabelled clinical EEG recordings, requiring only a small labelled probe set.
The Energy-Based Joint Embedding Predictive Architecture (EB-JEPA~\cite{ebjepa_lib})
provides a unified framework in which the SSL objective is expressed as an energy function
to be minimised, with anti-collapse regularisers preventing the trivial constant solution.

\paragraph{Report structure.}
\begin{description}[noitemsep]
  \item[Section~\ref{sec:sota}: State of the Art.]
    We first define three distinct \emph{scoring styles} used across the EEG literature,
    then provide the complete provenance table of published results on TUAB, verified
    against primary PDFs.
  \item[Section~\ref{sec:jepa}: EB-JEPA Model.]
    We describe the two-view energy concept, the dataset with its SSL/probe splits, and
    the full experimental results from all encoder $\times$ loss combinations tested.
  \item[Section~\ref{sec:math}: Per-Model Mathematical Descriptions.]
    For every model, we provide the complete energy function with formal component
    definitions, architecture details, the downstream MLP classifier, and the exact
    training/evaluation protocol.
  \item[Section~\ref{sec:best}: Combined Best Results.]
    We consolidate the best per-recording numbers from literature and this work.
  \item[Section~\ref{sec:conclusion}: Conclusion and Future Work.]
    Open scientific questions, methodological lessons, and prioritised future experiments.
\end{description}

\paragraph{Research question.}
Can a small encoder ($\leq$1.2M parameters) pre-trained without labels on 2\,717
EEG recordings learn representations that generalise to held-out patients and match
supervised models trained with labels?

%% ============================================================
\section{State of the Art on TUAB EEG Abnormality Detection}
\label{sec:sota}
%% ============================================================

\subsection{Scoring Styles}
\label{sec:scoring}

Published TUAB results use heterogeneous evaluation conventions.
We define four \emph{scoring styles} precisely:

\paragraph{(S1) Per-window (WIN).}
The evaluation set is segmented into all non-overlapping 10-second windows.
Each window is independently classified and inherits the binary label of its parent
recording.  BACC is computed over all individual window predictions.
For TUAB eval this yields ${\approx}$409\,k test windows from 2\,339 recordings
(the larger published eval split) or ${\approx}$400\,k from 276 recordings (our
official TUAB v3.0.1 split).
This is the convention of BIOT~\cite{biot}, LaBraM~\cite{labram}, ContraWR~\cite{contrawr},
and most published SSL models.
The test samples are \emph{not independent}: windows from the same recording share the
label and much of the signal, inflating effective sample size.

\paragraph{(S2) Per-recording via window embedding pooling (REC-WIN).}
A fixed number $N$ of evenly-spaced windows are extracted per recording and each is
encoded separately; the $N$ embeddings are mean-pooled into one vector.
One prediction is made per recording.
This is the default mode of our framework with $N{=}16$.
Noise averaging over 16 windows systematically raises BACC by ${\approx}5$\,pp over WIN
for the same encoder and probe.

\paragraph{(S3) Per-recording via majority vote (REC-VOTE).}
Each window is independently classified; the recording label is the majority class.
This style appears in some clinical pipelines but not in the SSL literature surveyed here.

\paragraph{(S4) Per-recording from the full recording (REC-FULL).}
The entire recording (variable length, up to 120 min) is fed to a model capable of
handling variable-length sequences (SSM, hierarchical model).
Used by FEMBA~\cite{femba} in its fine-tuned evaluation.

\medskip
\noindent\textbf{Critical warning.}
Comparing REC-WIN or REC-FULL numbers to WIN numbers as if they were the same quantity
produces a spurious ${\approx}5$\,pp advantage for the per-recording model.
Section~\ref{sec:conclusion} describes how this gap caused an initial erroneous claim.

\subsection{State-of-the-Art Provenance Table}
\label{sec:sota_table}

Table~\ref{tab:sota} compiles every TUAB result we could verify against a primary PDF.
Every literature number is backed by a specific table in a specific paper.
Entries in \colorbox{sota}{yellow} come from benchmark papers (EEG-Bench~\cite{eeg_bench},
FEMBA~\cite{femba}) rather than original method papers.
Entries in \colorbox{ours}{green} are our own results.

\begin{table}[htbp]
\centering
\caption{%
  Provenance table: verified results on TUAB EEG abnormality detection.
  WIN = per-window BACC; REC = per-recording (pooled).
  FT = fine-tuned; Frozen = frozen encoder + probe.
  All WIN numbers use the standard BIOT/LaBraM/FEMBA/EEGPT convention (409k windows).
  Our own runs use the official TUAB v3.0.1 eval split (276 recordings).
  $^{\dag}$Accuracy (not BACC); TUAB near-balanced (54\%/46\%), so acc $\approx$ BACC within 2\,pp.
  $^{\ddag}$Protocol unclear; BACC derived from reported sensitivity + specificity.
  Checks (\checkmark) indicate the value appears verbatim in the cited paper.
}
\label{tab:sota}
\scriptsize
\setlength{\tabcolsep}{3pt}
\resizebox{\textwidth}{!}{%
\begin{tabular}{L{1.5cm}L{3.8cm}C{1.1cm}C{1.1cm}C{1.1cm}L{2.8cm}L{1.3cm}L{0.7cm}}
\toprule
\textbf{Style} & \textbf{Model} & \textbf{BACC} & \textbf{AUROC} & \textbf{Params} &
\textbf{Source} & \textbf{Mode} & \textbf{Ver.}\\
\midrule
\multicolumn{8}{l}{\textit{Foundation models (large external pretraining)}} \\
WIN & LaBraM-Huge~\cite{labram} & 0.8258 & 0.9162 & 369M & Jiang 2024, Tbl 1 & FT & \checkmark \\
WIN & LaBraM-Large~\cite{labram} & 0.8226 & 0.9127 & 46M & Jiang 2024, Tbl 1 & FT & \checkmark \\
WIN & FEMBA-Huge~\cite{femba} & 0.8182 & 0.8921 & 386M & Tegon 2025, Tbl II & FT & \checkmark \\
WIN & FEMBA-Large~\cite{femba} & 0.8147 & 0.8856 & 77.8M & Tegon 2025, Tbl II & FT & \checkmark \\
WIN & LaBraM-Base~\cite{labram} & 0.8140 & 0.9022 & 5.8M & Jiang 2024, Tbl 1 & FT & \checkmark \\
WIN & FEMBA-Base~\cite{femba} & 0.8105 & 0.8829 & 47.7M & Tegon 2025, Tbl II & FT & \checkmark \\
WIN & EEG2Rep~\cite{femba} & 0.8052 & -- & -- & Tegon 2025, Tbl II & FT & \checkmark \\
\midrule
\multicolumn{8}{l}{\textit{Self-supervised models (TUAB-only or limited external pretraining)}} \\
WIN & BIOT~\cite{biot} (pretrained, PREST+SHHS) & 0.8019 & 0.8739 & 3.2M & Yang 2023, Tbl 4 & Frozen & \checkmark \\
WIN & AFTA & 0.8002 & 0.8848 & -- & Ref.\ not public & FT & \checkmark \\
WIN & EEGPT-25M~\cite{eegpt} & 0.7983 & 0.8718 & 25M & Wang 2024, Tbl 3+9 & Frozen & \checkmark \\
WIN & ST-Transformer~\cite{sttransformer} & 0.7966 & 0.8707 & 3.5M & LaBraM Tbl 1 / BIOT Tbl 4 & FT & \checkmark \\
WIN & BIOT~\cite{biot} (6 EEG datasets) & 0.7959 & 0.8815 & 3.2M & Yang 2023, Tbl 4 & Frozen & \checkmark \\
WIN & EEGPT-Tiny~\cite{eegpt} & 0.7959 & 0.8716 & 4.7M & Wang 2024, Tbl 3+9 & Frozen & \checkmark \\
WIN & BIOT~\cite{biot} (vanilla supervised) & 0.7925 & 0.8691 & 3.2M & Yang 2023, Tbl 4 & FT & \checkmark \\
WIN & SPaRCNet~\cite{sparcnet} & 0.7896 & 0.8676 & 0.79M & LaBraM Tbl 1 / BIOT Tbl 4 & FT & \checkmark \\
WIN & FFCL~\cite{ffcl} & 0.7848 & 0.8569 & 2.4M & Jiang 2024, Tbl 1 & FT & \checkmark \\
WIN & CNN-Transformer & 0.7777 & 0.8461 & 3.2M & Jiang 2024, Tbl 1 & FT & \checkmark \\
WIN & ContraWR~\cite{contrawr} & 0.7746 & 0.8456 & 1.6M & Jiang 2024, Tbl 1 & Frozen & \checkmark \\
WIN & BENDR~\cite{bendr} (FEMBA re-run) & 0.7696 & 0.8397 & 0.39M & Tegon 2025, Tbl II & FT & \checkmark \\
\midrule
\multicolumn{8}{l}{\textit{Supervised CNN baselines (Schirrmeister 2017 / spectral audit)}} \\
WIN & Deep4Net$^{\ddag}$ & 0.846 & -- & -- & Schirrmeister 2017, Tbl II & FT & \checkmark \\
WIN & ShallowConvNet$^{\ddag}$ & 0.839 & -- & 44k & Schirrmeister 2017, Tbl II & FT & \checkmark \\
WIN & EEGNet~\cite{eegnet} & 0.804 & -- & 9k & Spectral Audit~\cite{spectral_audit}, Tbl 1 & FT & \checkmark \\
WIN & Deep4Net~\cite{shallowconvnet} & 0.816 & -- & -- & Spectral Audit~\cite{spectral_audit}, Tbl 1 & FT & \checkmark \\
\midrule
\multicolumn{8}{l}{\textit{\cellcolor{ours}EB-JEPA (this work) --- WIN (per-window)}} \\
WIN & \cellcolor{ours}Conv1D + VICReg$^{*}$ (base) & \cellcolor{ours}0.756 & \cellcolor{ours}-- & \cellcolor{ours}0.4M & \cellcolor{ours}This work & \cellcolor{ours}Frozen & \\
WIN & \cellcolor{ours}Conv1D + VICReg$^{*}$ + corrupt & \cellcolor{ours}0.770 & \cellcolor{ours}-- & \cellcolor{ours}0.4M & \cellcolor{ours}This work & \cellcolor{ours}Frozen & \\
WIN & \cellcolor{ours}Conv1D + SIGReg + corrupt & \cellcolor{ours}\textbf{0.775} & \cellcolor{ours}-- & \cellcolor{ours}0.4M & \cellcolor{ours}This work & \cellcolor{ours}Frozen & \\
WIN & \cellcolor{ours}Conv1D + VICReg$^{*}$ + spec 0.1 & \cellcolor{ours}0.765 & \cellcolor{ours}-- & \cellcolor{ours}0.4M & \cellcolor{ours}This work & \cellcolor{ours}Frozen & \\
\midrule
\multicolumn{8}{l}{\textit{\cellcolor{ours}EB-JEPA (this work) --- REC-WIN ($N{=}16$, per-recording)}} \\
REC-WIN & \cellcolor{ours}Conv1D + VICReg (fixed) & \cellcolor{ours}0.819 & \cellcolor{ours}0.900 & \cellcolor{ours}0.4M & \cellcolor{ours}This work & \cellcolor{ours}Frozen & \\
REC-WIN & \cellcolor{ours}Conv1D + SIGReg + corrupt & \cellcolor{ours}0.825 & \cellcolor{ours}0.904 & \cellcolor{ours}0.4M & \cellcolor{ours}This work & \cellcolor{ours}Frozen & \\
REC-WIN & \cellcolor{ours}Conv1D + VICReg$^{*}$ + spec 0.1 & \cellcolor{ours}\textbf{0.836} & \cellcolor{ours}0.887 & \cellcolor{ours}0.4M & \cellcolor{ours}This work (buggy) & \cellcolor{ours}Frozen & \\
REC-WIN & \cellcolor{ours}Conv1D + VICReg$^{*}$ + corrupt (FT) & \cellcolor{ours}0.812 & \cellcolor{ours}0.908 & \cellcolor{ours}0.4M & \cellcolor{ours}This work & \cellcolor{ours}FT & \\
\bottomrule
\multicolumn{8}{l}{VICReg$^{*}$ = buggy \texttt{inv\_coeff=1}; VICReg (fixed) = paper-correct \texttt{inv\_coeff=25}.}
\end{tabular}}
\end{table}

\paragraph{Key observations.}
\begin{enumerate}[noitemsep]
  \item At WIN protocol, the state of the art is LaBraM-Huge (0.826) requiring 369M parameters
    pretrained on massive external data.  Among TUAB-only methods: LaBraM-Base (0.814) and
    BIOT-pretrained (0.802) are the targets.
  \item Our best WIN result (0.775, Conv1D + SIGReg + corrupt) matches ContraWR (0.775) and
    is 3.9\,pp below LaBraM-Base, with only 0.4M parameters and 20 training epochs.
  \item The REC-WIN vs WIN gap is $\approx$5\,pp for the same encoder: our best REC-WIN (0.836,
    buggy VICReg) vs WIN (0.765 same run) gives a 7.1\,pp gap.  This gap was initially used
    to claim we ``beat LaBraM,'' which is incorrect.
  \item The spectral audit~\cite{spectral_audit} finds that many published methods rely on
    the aperiodic (1/$f$) spectral slope for $\Delta$BACC 0.07--0.13; whether EB-JEPA representations
    share this shortcut is an open question.
\end{enumerate}

%% ============================================================
\section{EB-JEPA Model}
\label{sec:jepa}
%% ============================================================

\subsection{EB-JEPA Two-View Energy Concept}
\label{sec:jepa_concept}

EB-JEPA is a two-view energy-based SSL framework for EEG, inspired by the Joint Embedding
Predictive Architecture paradigm~\cite{ijepa,lecun2022}.  The key idea is:

\textbf{Starting from one EEG window, two corrupted views are created via independent
augmentation pathways (a) and (b), and an encoder (the ``world model'') is trained to
map both views to the same point in latent space.}

\paragraph{Formal setup.}
Given a 10-second EEG window $\mathbf{x} \in \mathbb{R}^{C \times T}$ with $C{=}19$ channels
and $T{=}2000$ samples at 200\,Hz:
\begin{itemize}[noitemsep]
  \item \textbf{Pathway (a)}: $\tilde{\mathbf{x}}_a = \mathcal{A}_a(\mathbf{x})$ --- apply
    augmentation stack $\mathcal{A}_a$ with independently sampled parameters.
  \item \textbf{Pathway (b)}: $\tilde{\mathbf{x}}_b = \mathcal{A}_b(\mathbf{x})$ --- apply
    the same augmentation family with a different random seed.
\end{itemize}
Both views are corrupted versions of the same brain state.  A shared encoder $f_\theta$
(the world model) and projector $g_\phi$ map each view to a latent vector:
\begin{equation}
  \mathbf{z}_a = g_\phi(f_\theta(\tilde{\mathbf{x}}_a)), \quad
  \mathbf{z}_b = g_\phi(f_\theta(\tilde{\mathbf{x}}_b)), \quad
  \mathbf{z}_a, \mathbf{z}_b \in \mathbb{R}^{d_p}, \quad d_p{=}1024.
\end{equation}
The energy function $E(\mathbf{z}_a, \mathbf{z}_b)$ is minimised during training:
\begin{equation}
  E = \underbrace{\|\mathbf{z}_a - \mathbf{z}_b\|^2}_{\text{invariance}}
    + w_V\, \underbrace{L_V(\mathbf{z}_a, \mathbf{z}_b)}_{\text{anti-collapse (variance)}}
    + w_C\, \underbrace{L_C(\mathbf{z}_a, \mathbf{z}_b)}_{\text{decorrelation (covariance)}}.
\end{equation}
Without the anti-collapse terms, the trivial global minimum $\mathbf{z}_a = \mathbf{z}_b = \mathbf{0}$
satisfies $E{=}0$ while encoding no information.  VICReg~\cite{vicreg} and SIGReg prevent this.

The encoder acts as a \textbf{world model}: it abstracts away the specific corruption and
retains the clinically relevant brain state.  Pathological brains manifest statistically
different signal statistics across all augmentation families; the encoder is forced to
represent what is preserved across corruptions, which correlates with pathology.

\subsection{EEG Dataset: SSL Training vs.\ Probe Evaluation Splits}
\label{sec:dataset}

\paragraph{Dataset.}
All experiments use the \textbf{TUH Abnormal EEG Corpus (TUAB) v3.0.1}~\cite{tuab}.
19-channel scalp EEG in the 10-20 system, binary normal/abnormal labels assigned
by a board-certified neurologist per recording.

\paragraph{Offline preprocessing (provided by dataset).}
Applied once to the raw files before any experiment:
\begin{enumerate}[noitemsep]
  \item \textbf{Channel selection}: 19 standard scalp electrodes (FP1, FP2, F3, F4, C3, C4,
    P3, P4, O1, O2, F7, F8, T3, T4, T5, T6, FZ, CZ, PZ).
  \item \textbf{Resampling to 200\,Hz}: raw rates vary (250/256/512\,Hz); resampled to a
    common grid so that a 10-second window $= 2000$ samples, convenient for stride-2 convolutions.
  \item \textbf{60\,Hz notch filter}: US mains frequency; removed before learning.
  \item \textbf{0.1--75\,Hz bandpass}: standard clinical EEG band covering delta through low gamma;
    discards DC drift and high-frequency artefacts above the 100\,Hz Nyquist.
\end{enumerate}

\paragraph{Online preprocessing (per-window, at load time).}
Each 10-second window $\mathbf{x}_c \in \mathbb{R}^T$ per channel $c$ is z-scored:
\begin{equation}
  \hat{x}_c[t] = \frac{x_c[t] - \mu_c}{\sigma_c + \varepsilon}, \quad
  \mu_c = \frac{1}{T}\sum_t x_c[t], \quad
  \sigma_c = \sqrt{\frac{1}{T}\sum_t (x_c[t]-\mu_c)^2}, \quad
  \varepsilon = 10^{-6}.
\end{equation}
This is computed \emph{within the window only} (no cross-window statistics), removing
per-electrode amplitude offsets caused by impedance variability.

\paragraph{Patient-disjoint splits.}
\begin{itemize}[noitemsep]
  \item \textbf{SSL training split} (2\,717 recordings: 1\,385 normal, 1\,332 abnormal):
    used for SSL pre-training with labels discarded.  The encoder sees only raw EEG.
  \item \textbf{Evaluation split} (276 recordings: 150 normal, 126 abnormal):
    used for final performance estimation only.  No hyperparameter selection on this split.
\end{itemize}
No patient appears in both splits.

\paragraph{Probe training (ssl split with labels).}
After SSL pre-training, the encoder is frozen.  An MLP probe is trained on the same
2\,717 recordings using their normal/abnormal labels.  For each recording, $N{=}16$
evenly-spaced 10-second windows are extracted, encoded, and mean-pooled into one
256-dim vector (REC-WIN protocol).  The probe trains on these pooled vectors.

\paragraph{Evaluation.}
On the 276-recording eval split, both WIN and REC-WIN metrics are computed.
No test-time augmentation, no threshold tuning on the eval split.

\paragraph{Transfer evaluation on TUEV.}
Selected encoders are also evaluated on the \textbf{TUH EEG Events corpus (TUEV v2.0.0)},
a 6-class event detection task (spike-and-slow-wave, generalised periodic discharge,
lateralised periodic discharge, eye movement, artefact, background).
A frozen logistic regression probe is trained on event-centred 10-second windows
from TUEV train, evaluated on TUEV eval.  No TUEV data was used during SSL pretraining;
this is a pure transfer benchmark.

For context, Table~\ref{tab:tuev_supervised} reports supervised baselines trained
\emph{end-to-end on TUEV} using BIOT's official pipeline (AdamW, lr=1e-3, 10 epochs).
Our SSL encoder (pretrained on TUAB, no TUEV labels) plus a linear probe competes
with several of these fully supervised models.

\begin{table}[h]
\centering
\caption{Supervised TUEV baselines (our re-run, BIOT pipeline, end-to-end on TUEV).
  Compare to our SSL transfer: best 0.425 (confounded) / clean best 0.389 (NaiveJEPA).}
\label{tab:tuev_supervised}
\small
\begin{tabular}{lccc}
\toprule
\textbf{Model} & \textbf{Params} & \textbf{Our TUEV BACC} & \textbf{Lit.\ TUEV BACC} \\
\midrule
BIOT pretrained (PREST-16ch) & 3.2M & 0.5009 & 0.5281$^{\dagger}$ \\
BIOT vanilla (supervised) & 3.2M & 0.4994 & 0.4682 \\
SPaRCNet & 0.79M & 0.4726 & 0.4161 \\
ContraWR & 1.6M & 0.4664 & 0.4384 \\
FFCL & 2.4M & 0.4551 & 0.3979 \\
CNN-Transformer & 3.2M & 0.4100 & 0.4087 \\
ST-Transformer & 3.5M & 0.3385 & 0.3984 \\
\midrule
\textbf{Ours: NaiveJEPA} (TUAB SSL $\to$ TUEV probe) & 0.4M & \textbf{0.389} & --- \\
\textbf{Ours: VICReg} ($\lambda{=}1$, TUAB SSL $\to$ probe) & 0.4M & 0.382 & --- \\
\bottomrule
\multicolumn{4}{l}{$^{\dagger}$Literature uses 18-ch EEG-six-datasets ckpt; we used 16-ch PREST ckpt.}
\end{tabular}
\end{table}

Our SSL transfer (0.382--0.389 BACC) beats ST-Transformer (0.3385) and CNN-Transformer (0.4100)
despite never training on TUEV labels, using only 0.4M parameters and a linear probe.
The gap to supervised BIOT (0.50) represents the headroom of full task-specific training.

\subsection{Full Results Table}
\label{sec:full_results}

Table~\ref{tab:full_results} reports all encoder $\times$ loss $\times$ augmentation
combinations evaluated.  The TUEV column quantifies transfer generalisation.

\begin{table}[htbp]
\centering
\caption{%
  Complete EB-JEPA results on TUAB v3.0.1 (276-recording eval split) and TUEV v2.0.0
  (6-class transfer).
  WIN = per-window BACC; REC = per-recording BACC ($N{=}16$ mean-pooled); AUROC = area under ROC.
  BACC\textsubscript{TUEV} = 6-class macro BACC (frozen logistic probe, event-centred windows).
  VICReg$^{*}$ = $\lambda_{\rm inv}{=}1$ (default, 3\,pp better than paper-correct $\lambda{=}25$);
  VICReg = $\lambda_{\rm inv}{=}25$ (paper-correct, worse on this EEG task).
  ``--'' = not evaluated.
  $^{a}$Accuracy (not BACC), early eval script.
}
\label{tab:full_results}
\small
\setlength{\tabcolsep}{3pt}
\resizebox{\textwidth}{!}{%
\begin{tabular}{L{4.4cm}C{1.2cm}C{1.2cm}C{1.2cm}C{1.2cm}L{1.6cm}L{1.2cm}L{1.2cm}}
\toprule
\textbf{Model / Config} &
\textbf{BACC\textsubscript{WIN}} &
\textbf{BACC\textsubscript{REC}} &
\textbf{AUROC\textsubscript{REC}} &
\textbf{BACC\textsubscript{TUEV}} &
\textbf{Score style} &
\textbf{Probe} &
\textbf{Data} \\
\midrule
\multicolumn{8}{l}{\textit{Conv1D (0.4M) --- spectral coeff sweep, buggy VICReg$^{*}$}} \\
Conv + VICReg$^{*}$ (base) & 0.756 & 0.796 & 0.888 & 0.381 & WIN / REC-WIN & Frozen & TUAB \\
Conv + VICReg$^{*}$ + spec 0.05 & -- & 0.821 & 0.884 & -- & REC-WIN & Frozen & TUAB \\
Conv + VICReg$^{*}$ + spec 0.1 & 0.765 & 0.836 & 0.887 & \textbf{0.425} & WIN / REC-WIN & Frozen & TUAB \\
Conv + VICReg$^{*}$ + spec 0.3 & -- & 0.789 & 0.877 & -- & REC-WIN & Frozen & TUAB \\
Conv + VICReg$^{*}$ + spec 1.0 & -- & 0.785 & 0.874 & -- & REC-WIN & Frozen & TUAB \\
Conv + VICReg$^{*}$ + fftc 0.05 & -- & 0.807 & 0.887 & -- & REC-WIN & Frozen & TUAB \\
Conv + VICReg$^{*}$ + fftc 0.1 & -- & 0.798 & 0.881 & -- & REC-WIN & Frozen & TUAB \\
Conv + VICReg$^{*}$ + fftc 0.3 & -- & 0.802 & 0.889 & -- & REC-WIN & Frozen & TUAB \\
\midrule
\multicolumn{8}{l}{\textit{Conv1D --- spectral coeff sweep, FIXED VICReg (this work)}} \\
Conv + VICReg (fixed, spec=0) & -- & 0.819 & 0.900 & 0.382 & REC-WIN & Frozen & TUAB \\
Conv + VICReg (fixed) + spec 0.05 & -- & 0.802 & 0.891 & 0.313 & REC-WIN & Frozen & TUAB \\
Conv + VICReg (fixed) + spec 0.10 & -- & 0.795 & 0.892 & 0.375 & REC-WIN & Frozen & TUAB \\
Conv + VICReg (fixed) + spec 0.20 & -- & 0.813 & 0.889 & 0.357 & REC-WIN & Frozen & TUAB \\
Conv + VICReg (fixed) + spec 0.30 & -- & 0.799 & 0.885 & 0.375 & REC-WIN & Frozen & TUAB \\
\midrule
\multicolumn{8}{l}{\textit{Conv1D --- corruption augmentation ablation (3 seeds)}} \\
Conv + VICReg$^{*}$ + corrupt (s0) & 0.770 & 0.825 & 0.904 & 0.382 & WIN / REC-WIN & Frozen & TUAB \\
Conv + VICReg$^{*}$ + corrupt (s1) & -- & 0.816 & 0.891 & -- & REC-WIN & Frozen & TUAB \\
Conv + VICReg$^{*}$ + corrupt (s2) & -- & 0.818 & 0.906 & -- & REC-WIN & Frozen & TUAB \\
Conv + corrupt + spec 0.1 & -- & 0.802 & 0.899 & -- & REC-WIN & Frozen & TUAB \\
Conv + SIGReg + corrupt & 0.775 & 0.825 & 0.904 & 0.363 & WIN / REC-WIN & Frozen & TUAB \\
\midrule
\multicolumn{8}{l}{\textit{Conv1D --- fine-tuning, multi-corpus, and new objectives}} \\
Conv + VICReg$^{*}$ + corrupt (FT) & -- & 0.812 & 0.908 & -- & REC-WIN & Fine-tuned & TUAB \\
Conv + VICReg$^{*}$ + multicorpus & -- & 0.812 & -- & 0.394 & REC-WIN & Frozen & TUAB+TUEV+TUSZ+TUEP \\
Conv + Pred-JEPA (20ep) & -- & 0.808 & 0.899 & -- & REC-WIN & Frozen & TUAB \\
Conv + NaiveJEPA (20ep) & -- & 0.779 & 0.863 & 0.389 & REC-WIN & Frozen & TUAB \\
EMA-JEPA on Conv (150ep) & -- & 0.764 & -- & -- & REC-WIN & Frozen & TUAB \\
\midrule
\multicolumn{8}{l}{\textit{LaBraM-style patch Transformer (1.2M)}} \\
LaBraM + VICReg$^{*}$ & 0.775$^{a}$ & -- & -- & -- & WIN & Frozen & TUAB \\
LaBraM + SIGReg & 0.725$^{a}$ & -- & -- & -- & WIN & Frozen & TUAB \\
LaBraM + Pred-JEPA (20ep) & -- & 0.647 & -- & -- & REC-WIN & Frozen & TUAB \\
\midrule
\multicolumn{8}{l}{\textit{BIOT-style FFT Transformer (1.2M)}} \\
BIOT + VICReg$^{*}$ & 0.775$^{a}$ & -- & -- & -- & WIN & Frozen & TUAB \\
BIOT + SIGReg & 0.783$^{a}$ & -- & -- & -- & WIN & Frozen & TUAB \\
BIOT + Pred-JEPA (20ep) & -- & 0.790 & 0.866 & -- & REC-WIN & Frozen & TUAB \\
\midrule
\multicolumn{8}{l}{\textit{EEGPT-style channel-pool Transformer (0.8M)}} \\
EEGPT + VICReg$^{*}$ & collapsed & -- & -- & -- & WIN & Frozen & TUAB \\
EEGPT + SIGReg & collapsed & -- & -- & -- & WIN & Frozen & TUAB \\
EEGPT + Pred-JEPA (20ep) & -- & 0.613 & 0.642 & -- & REC-WIN & Frozen & TUAB \\
EEGPT + VICReg (fixed) & -- & -- & -- & 0.182 & -- & Frozen & TUAB \\
EEGPT + VICReg (no chan drop) & -- & -- & -- & 0.358 & -- & Frozen & TUAB \\
\bottomrule
\end{tabular}}
\end{table}

\paragraph{Key findings.}
\begin{enumerate}[noitemsep]
  \item \textbf{Conv1D dominates} all encoder families at equal or lower parameter count.
  \item \textbf{Corruption augmentation} (time masking + spike injection) is the single
    most impactful ingredient: $+0.014$ WIN, $+0.029$ REC.
  \item \textbf{Spectral loss is neutral with fixed VICReg}: a full sweep of
    $w_{\rm spec} \in \{0.05, 0.10, 0.20, 0.30\}$ with correct VICReg coefficients (inv=25)
    yields TUAB BACC in $[0.795, 0.813]$, all within single-seed noise of the no-spectral
    baseline (0.819, $\pm$0.004 across 3 seeds).  Drops range from 0.006 to 0.024\,pp ---
    not statistically significant with single-seed runs.
    The previous best of 0.836 (spec=0.1, buggy VICReg) is not reproducible with correct
    coefficients; the spectral loss was compensating for the 25$\times$ underweighted
    invariance term, not providing independent signal.  The accurate statement is
    \emph{neutral-to-mildly-harmful}, not dramatically bad.
  \item \textbf{EEGPT collapsed} under both objectives with the buggy VICReg;
    the channel-drop vs.\ channel-pool conflict is the architectural root cause.
    Removing channel-drop improves TUEV transfer (0.182 $\to$ 0.358).
  \item \textbf{NaiveJEPA} (temporal split, no anti-collapse term) reaches 0.779 BACC
    and 0.389 TUEV, competitive with VICReg at a fraction of the design complexity.
  \item \textbf{TUEV best}: Conv + VICReg$^{*}$ + spec 0.1 = 0.425 (buggy VICReg run).
    This cross-task transfer result is the strongest evidence of general EEG representation
    quality, as TUEV labels were never seen during pretraining.
\end{enumerate}

%% ============================================================
\section{Per-Model Mathematical Descriptions}
\label{sec:math}
%% ============================================================

\subsection{Conv1D Encoder + VICReg}
\label{sec:math_conv}

\paragraph{Energy function.}
\begin{equation}
  \boxed{E_{\mathrm{VICReg}} = w_2 \cdot L_2 + w_V \cdot L_V + w_C \cdot L_C}
  \quad (w_2, w_V, w_C) = (25,\,25,\,1).
\end{equation}

\textbf{$L_2$ (invariance):}
\begin{equation}
  L_2(\mathbf{Z}, \mathbf{Z}') = \frac{1}{N}\sum_{i=1}^{N}\|\mathbf{z}_{a,i} - \mathbf{z}_{b,i}\|^2,
  \quad \mathbf{z}_{a,i}, \mathbf{z}_{b,i} \in \mathbb{R}^{d_p},\; d_p{=}1024.
\end{equation}
This is the core energy: it drives both views toward the same latent point.

\textbf{$L_V$ (variance / anti-collapse hinge):}
\begin{equation}
  L_V = \frac{1}{2d_p}\sum_{j=1}^{d_p}\Bigl[
    \max(0,\,\gamma - \mathrm{Std}(\mathbf{Z}_{a,:,j})) +
    \max(0,\,\gamma - \mathrm{Std}(\mathbf{Z}_{b,:,j}))
  \Bigr], \quad \gamma = 1.
\end{equation}
$\mathbf{Z}_a \in \mathbb{R}^{N \times d_p}$ is the batch of projected embeddings for view $a$.
$L_V$ penalises any dimension $j$ whose within-batch standard deviation falls below 1,
preventing collapse to a constant vector.

\textbf{$L_C$ (covariance / decorrelation):}
\begin{equation}
  L_C = \frac{1}{2d_p(d_p{-}1)}\Bigl[
    \sum_{i\neq j}\mathrm{Cov}(\mathbf{Z}_a)^2_{ij} +
    \sum_{i\neq j}\mathrm{Cov}(\mathbf{Z}_b)^2_{ij}
  \Bigr],
\end{equation}
where $\mathrm{Cov}(\mathbf{Z}) = (\mathbf{Z}-\bar{\mathbf{Z}})^\top(\mathbf{Z}-\bar{\mathbf{Z}})/(N-1)$.
$L_C$ penalises correlated dimensions, encouraging the encoder to spread information
across all $d_p$ independent directions.

\paragraph{Architecture.}
The Conv1D encoder is designed from scratch.  Input $\mathbf{X} \in \mathbb{R}^{19 \times 2000}$:
\begin{equation}
  \mathbf{H}^{(l)} = \mathrm{GELU}\!\Bigl(\mathrm{BN}\bigl(
    \mathrm{Conv1d}_{F_l,\,k=7,\,s=2}(\mathbf{H}^{(l-1)})\bigr)\Bigr),
  \quad l=1,\ldots,4,
\end{equation}
with channel widths $F_0{=}19, F_1{=}64, F_2{=}128, F_3{=}256, F_4{=}256$.
Stride 2 halves the time axis each block: $2000 \to 1000 \to 500 \to 250 \to 125$.
Global average pooling over time, then $\mathbf{W}_{\rm out} \in \mathbb{R}^{256\times256}$:
\begin{equation}
  \mathbf{z} = \mathbf{W}_{\rm out}\!\left(\tfrac{1}{125}\textstyle\sum_{t=1}^{125}
    \mathbf{H}^{(4)}_{:,t}\right)\in\mathbb{R}^{256}.
\end{equation}
Projector MLP: $256\to1024\to1024$ (two linear layers with batch norm and ReLU).
Total parameters: $\approx 0.4$M.

\subsection{Conv1D + SIGReg}
\label{sec:math_sigreg}

\paragraph{Energy function.}
SIGReg replaces $L_V$ and $L_C$ with a Gaussianisation surrogate~\cite{epps1983}:
\begin{equation}
  \boxed{E_{\mathrm{SIGReg}} = L_2 + \lambda \cdot L_{\mathrm{SIG}}},
  \quad \lambda = 10,\; K{=}256.
\end{equation}
For $K$ random unit-sphere directions $\mathbf{v}_k \sim \mathcal{U}(\mathcal{S}^{d_p-1})$:
\begin{equation}
  L_{\mathrm{SIG}} = \frac{1}{2K}\sum_{k=1}^{K}\Bigl[
    \mathrm{EP}_N\!\bigl(\{\mathbf{v}_k^\top \mathbf{z}_{a,i}\}_i\bigr) +
    \mathrm{EP}_N\!\bigl(\{\mathbf{v}_k^\top \mathbf{z}_{b,i}\}_i\bigr)\Bigr],
\end{equation}
where $\mathrm{EP}_N(\{u_i\})$ is the Epps--Pulley test statistic measuring deviation
of the empirical distribution from Gaussian:
\begin{equation}
  \mathrm{EP}_N(\{u_i\}) = \int_{-3}^{3}
    e^{-t^2/2}\left|\hat\varphi_N(t) - e^{-t^2/2}\right|^2 dt,
  \quad
  \hat\varphi_N(t) = \tfrac{1}{N}\textstyle\sum_{i=1}^N e^{jtu_i}.
\end{equation}
The integral is approximated at 10 equally-spaced points in $[-3,3]$.
Collapse (constant vectors) forces projections to be non-Gaussian, which the EP test
detects.  This is a strictly stronger collapse criterion than VICReg's variance hinge.

\paragraph{Architecture.} Identical to Conv1D + VICReg above.

\subsection{Conv1D + Spectral Regularisation}
\label{sec:math_spectral}

\paragraph{Energy function.}
\begin{equation}
  \boxed{E_{\mathrm{spec}} = E_{\mathrm{VICReg}} + w_{\rm spec} \cdot L_{\rm spec}},
  \quad w_{\rm spec} \in \{0.05, 0.1, 0.3, 1.0\}.
\end{equation}
$L_{\rm spec}$ is a DDSP-inspired multi-scale spectral loss~\cite{engel2020ddsp} applied
to the encoder \emph{feature maps} (not the raw signal) of both views:
\begin{equation}
  L_{\rm spec} = \frac{1}{|M|}\sum_{m \in M} \frac{1}{N f_m}\sum_{i=1}^N\sum_{k=1}^{f_m}
    \Bigl(\log|\hat{\mathbf{S}}^{(m)}_{i,k}| - \log|\mathbf{S}^{(m)}_{i,k}|\Bigr)^2,
\end{equation}
where $M = \{256, 128, 64, 32\}$ are the STFT window sizes, $f_m = m/2 + 1$ is the
number of frequency bins, $\mathbf{S}^{(m)}_{i,k}$ is the STFT magnitude of the
\emph{target} feature map at scale $m$, and $\hat{\mathbf{S}}^{(m)}_{i,k}$ is that of
the \emph{context} feature map.  This encourages the encoder to map views with similar
spectral content to similar feature maps.

\paragraph{Key finding from this work.}
The spectral loss helped with the \emph{buggy} VICReg (inv=1) but hurts with the
\emph{fixed} VICReg (inv=25).  With correct invariance weighting, VICReg alone dominates
and the spectral term conflicts rather than supplements it.

\subsection{LaBraM-Style Patch Transformer}
\label{sec:math_labram}

\paragraph{Energy function.}
$E_{\mathrm{LaBraM}} = 25 L_2 + 25 L_V + L_C$ (VICReg) or $L_2 + 10 L_{\mathrm{SIG}}$
(SIGReg).  Identical form to Conv1D; architecture differs.

\paragraph{Architecture.}
Inspired by LaBraM~\cite{labram} but built from scratch (without the VQ-codebook
neural tokeniser).  Each channel is split into $P{=}10$ non-overlapping patches of
$w{=}200$ samples each:
\begin{equation}
  \mathbf{p}_{c,k} = \mathbf{X}_{c,\,(k-1)w:kw}\in\mathbb{R}^{w},\quad c\in[C],\;k\in[P].
\end{equation}
Linear patch embedding plus learnable temporal and channel embeddings:
\begin{equation}
  \hat{\mathbf{e}}_{c,k} = \mathbf{W}_p\mathbf{p}_{c,k} + \mathbf{te}_k + \mathbf{se}_c
  \in\mathbb{R}^{d_e},\quad d_e{=}128.
\end{equation}
Sequence of $C \times P = 190$ tokens fed to 4 Transformer blocks~\cite{transformer}
(multi-head attention with $n_h{=}4$ heads, pre-LN, FFN ratio 4):
\begin{equation}
  \mathrm{Attn}(Q,K,V) = \mathrm{softmax}\!\left(\frac{QK^\top}{\sqrt{d_h}}\right)V,
  \quad d_h = d_e / n_h = 32.
\end{equation}
Mean-pool over 190 tokens; linear projection $d_e\to 256$: $\mathbf{z}\in\mathbb{R}^{256}$.
Total parameters: $\approx 1.2$M.

\subsection{BIOT-Style FFT Transformer}
\label{sec:math_biot}

\paragraph{Energy function.} Same as LaBraM-style.

\paragraph{Architecture.}
Inspired by BIOT~\cite{biot} (from scratch).  The key difference: each patch
$\mathbf{p}_{c,k}$ is embedded via its one-sided FFT magnitude spectrum:
\begin{equation}
  \mathbf{s}_{c,k}[m] = \left|\sum_{n=0}^{w-1}\mathbf{p}_{c,k}[n]\,
    e^{-j2\pi mn/w}\right|,\quad m=0,\ldots,\tfrac{w}{2}.
\end{equation}
This yields a $(w/2+1){=}101$-dimensional frequency vector per patch.
Linear projection to $d_e{=}128$, then the same Transformer and pooling as LaBraM-style.
Total parameters: $\approx 1.2$M.

\paragraph{Failure mode.}
The FFT discards phase information.  EEG pathology often manifests in temporal dynamics
(spike-wave complexes, burst-suppression onset) that depend on phase; discarding it limits
BIOT-style to frequency-content features.

\subsection{EEGPT-Style Hierarchical Channel-Pool Transformer}
\label{sec:math_eegpt}

\paragraph{Energy function.} Same as LaBraM-style.

\paragraph{Architecture.}
Inspired by EEGPT~\cite{eegpt} (from scratch).  A learnable \emph{channel-pooling} step
reduces 19 channel tokens to 1 summary token per time patch:

Each patch is linearly embedded with a learnable channel codebook:
$\mathrm{tok}_{c,k} = \mathbf{W}_p\mathbf{p}_{c,k} + \mathbf{s}_c \in \mathbb{R}^{d_e}$.
A global query $\mathbf{q}\in\mathbb{R}^{d_e}$ attends over all $C{=}19$ channel tokens:
\begin{equation}
  \mathbf{g}_k = \mathrm{Attn}\!\bigl(\mathbf{q};\;
    [\mathrm{tok}_{1,k},\ldots,\mathrm{tok}_{C,k}]\bigr)\in\mathbb{R}^{d_e}.
\end{equation}
This reduces 190 tokens to 10 summary tokens.  A 4-block Transformer over the 10 tokens,
then mean-pool and linear projection to $\mathbb{R}^{256}$.
Total parameters: $\approx 0.8$M.

\paragraph{Failure mode: channel-pool $\times$ channel-drop conflict.}
The channel-drop augmentation ($p{=}0.2$) zeros different channels in view $a$ and $b$
independently.  The pooling attention $\mathbf{g}_k$ therefore receives a different
subset of channels per view, making the summary tokens structurally incompatible.
The invariance loss $L_2 = \|\mathbf{z}_a - \mathbf{z}_b\|^2$ then demands the encoder
ignore which channels are present, defeating the spatial-attention design.
This is why EEGPT collapsed with channel-drop augmentation; removing channel-drop
raises TUEV BACC from 0.182 to 0.358.

\subsection{NaiveJEPA (Temporal-Split Predictive)}
\label{sec:math_naive}

\paragraph{Concept.}
NaiveJEPA follows I-JEPA~\cite{ijepa} applied to EEG: a single window $\mathbf{x}$
is split in time at ratio $r$ (default 0.5).  An online encoder sees the first $r T$
samples (context); an EMA target encoder sees the remaining $(1-r)T$ samples (target).
A predictor maps context embedding to target embedding.  Anti-collapse is \emph{not needed}:
context and target cover disjoint time spans, so constant predictions cannot minimise MSE.

\paragraph{Energy function.}
\begin{equation}
  \boxed{E_{\mathrm{NaiveJEPA}} = L_{\mathrm{pred}} = \mathrm{MSE}(\hat{\mathbf{z}}_{\rm tgt},\;
    \mathbf{z}_{\rm tgt})},
\end{equation}
\begin{equation}
  \hat{\mathbf{z}}_{\rm tgt} = h_\psi(\bar{\mathbf{z}}_{\rm ctx}), \quad
  \bar{\mathbf{z}}_{\rm ctx} = \frac{1}{F_{\rm ctx}}\sum_{t=1}^{F_{\rm ctx}}
    f_\theta(\mathbf{x}_{:,:rT})_t, \quad
  \mathbf{z}_{\rm tgt} = \mathrm{sg}\!\left(
    \frac{1}{F_{\rm tgt}}\sum_{t=1}^{F_{\rm tgt}}
    f_{\bar\theta}(\mathbf{x}_{:,rT:})_t\right),
\end{equation}
where $f_\theta$ is the online (trained) encoder, $f_{\bar\theta}$ is the EMA target
encoder with no gradient, $h_\psi$ is a 2-layer MLP predictor
($D\to 2D\to D$ with LayerNorm and GELU), and $\mathrm{sg}(\cdot)$ is the stop-gradient.
The EMA update:
\begin{equation}
  \bar\theta \leftarrow m\,\bar\theta + (1-m)\,\theta, \quad m = 0.996.
\end{equation}

\paragraph{Architecture.} Same Conv1D encoder as above.  Predictor: LayerNorm + Linear
($D \to 2D$) + GELU + Linear ($2D \to D$), $D{=}256$.

\subsection{Predictive EMA-JEPA (GRU-based)}
\label{sec:math_pred}

\paragraph{Energy function.}
\begin{equation}
  \boxed{E_{\mathrm{Pred}} = L_{\rm pred} + w_V\,L_V^{\rm flat} + w_C\,L_C^{\rm flat}},
  \quad w_V = w_C = 1.
\end{equation}
Two independently augmented views $\mathbf{v}_1, \mathbf{v}_2$ of the \emph{same} window
are encoded to frame sequences $\mathbf{h}^{(1)}, \mathbf{h}^{(2)} \in \mathbb{R}^{F\times D}$.
An EMA target encoder provides $\mathbf{h}^{(2)}$ without gradient.
A GRU predictor rolls forward from frame 0 of $\mathbf{h}^{(1)}$:
\begin{equation}
  \hat{\mathbf{h}}^{(2)}_t = \mathrm{GRU}(\hat{\mathbf{h}}^{(2)}_{t-1}, \mathbf{0}),
  \quad
  L_{\rm pred} = \frac{1}{F-1}\sum_{t=1}^{F-1}
    \|\hat{\mathbf{h}}^{(2)}_t - \mathbf{h}^{(2)}_t\|^2.
\end{equation}
Anti-collapse: Hinge-std and Covariance losses applied to the flattened frame matrix
$\mathbf{H}^{(1)}\in\mathbb{R}^{(NF)\times D}$ (same $L_V$, $L_C$ definitions as VICReg).

\paragraph{Failure mode.}
Both views span the \emph{same} 10-second window; the GRU's ``temporal prediction'' is
actually augmentation-invariance in disguise.  The anti-collapse term (std\_loss) was
near-zero throughout training (encoder found near-constant solutions), and downstream
TUAB BACC (0.764--0.808) is below VICReg, confirming collapse.

\subsection{MLP Classifier (Probe)}
\label{sec:math_mlp}

The downstream classification probe is implemented as a 3-layer MLP:
\begin{align}
  \mathbf{h}_1 &= \mathrm{ReLU}(\mathbf{W}_1\mathbf{z} + \mathbf{b}_1),\quad \mathbf{W}_1\in\mathbb{R}^{128\times256},\\
  \mathbf{h}_2 &= \mathrm{ReLU}(\mathbf{W}_2\mathbf{h}_1 + \mathbf{b}_2),\quad \mathbf{W}_2\in\mathbb{R}^{64\times128},\\
  \hat{y} &= \mathrm{Softmax}(\mathbf{W}_3\mathbf{h}_2 + \mathbf{b}_3),\quad \mathbf{W}_3\in\mathbb{R}^{2\times64}.
\end{align}
Implemented via \texttt{sklearn.neural\_network.MLPClassifier}
(\texttt{hidden\_layer\_sizes=(128,64)}, \texttt{max\_iter=500}, \texttt{early\_stopping=True}).
A \texttt{StandardScaler} is fitted on training embeddings and applied identically at eval.

\textbf{Frozen mode}: only $\{\mathbf{W}_l,\mathbf{b}_l\}$ are updated; $f_\theta$ parameters
receive no gradient.\\
\textbf{Fine-tuned mode}: all parameters (encoder + probe) are optimised jointly at
learning rate $10^{-5}$ for the encoder, $10^{-3}$ for the probe.

Note: A logistic regression probe (sklearn LogReg) is the standard SSL community probe
and would allow direct comparison to SimCLR/MoCo literature.  Our MLP is more expressive
and may inflate absolute numbers by $\approx$0.01--0.02 BACC relative to LogReg.

\subsection{Training and Testing Protocol}
\label{sec:protocol}

\paragraph{SSL pre-training.}
\begin{itemize}[noitemsep]
  \item \textbf{Dataset}: TUAB SSL split (2\,717 recordings, labels discarded).
  \item \textbf{Windows}: non-overlapping 10-second windows, per-channel z-score online.
  \item \textbf{Optimiser}: AdamW, lr $= 10^{-3}$, weight decay $10^{-6}$.
  \item \textbf{Batch size}: 128 windows.
  \item \textbf{Duration}: 20 epochs.
  \item \textbf{Augmentation}: Gaussian noise ($\sigma{=}0.1$), amplitude jitter ($\pm$20\%),
    channel dropout ($p{=}0.2$), time masking (20\% of timepoints), spike injection ($\pm6\sigma$
    at 20\% timepoints) for the corruption variant.
\end{itemize}

\paragraph{Probe training (ssl split with labels).}
$N{=}16$ evenly-spaced windows per training recording $\to$ frozen encoder
$\to$ 16 embeddings $\to$ mean-pool $\to$ one 256-dim vector per recording.
Probe trained on 2\,717 vectors with their binary labels.

\paragraph{Evaluation (eval split).}
\begin{itemize}[noitemsep]
  \item \textbf{WIN}: all non-overlapping 10-second windows, one prediction each,
    BACC over all window predictions.
  \item \textbf{REC-WIN}: same $N{=}16$ extraction as probe training, mean-pool, classify.
  \item No test-time augmentation, no threshold tuning.
\end{itemize}

\paragraph{TUEV transfer probe.}
Frozen logistic regression on event-centred 10-second windows (5 events per class subsampled,
background subsampled to $\approx$1.26\% of windows), evaluated on the TUEV v2.0.0 eval split.
Window length: 10 seconds (2000 samples) for all encoders except Conv1D (5 seconds for early runs).

%% ============================================================
\section{Combined Best Results (Per-Recording Focus)}
\label{sec:best}
%% ============================================================

Table~\ref{tab:best} consolidates the best per-recording results from literature and this work,
filtered to REC-WIN or REC-FULL (one prediction per patient, the clinically meaningful unit).
Per-window results are included for reference where no REC protocol exists.

\begin{table}[htbp]
\centering
\caption{%
  Best TUAB results, per-recording focus.
  Bold = best in each block.
  FEMBA and LaBraM entries use massive external pretraining (TUEG 21kh); these are
  not directly comparable to TUAB-only methods.
  $^{\ddag}$Best REC of EEGNet reported in EEGNet literature;
  BACC derived by us from per-window sensitivity/specificity numbers.
}
\label{tab:best}
\small
\setlength{\tabcolsep}{4pt}
\resizebox{\textwidth}{!}{%
\begin{tabular}{llC{1.1cm}C{1.3cm}C{1.3cm}C{1.3cm}L{1.4cm}}
\toprule
\textbf{Method} & \textbf{Encoder} & \textbf{Params} &
  \textbf{BACC\textsubscript{WIN}} &
  \textbf{BACC\textsubscript{REC}} &
  \textbf{AUROC\textsubscript{REC}} &
  \textbf{Mode} \\
\midrule
\multicolumn{7}{l}{\textit{Foundation models with large external pretraining}} \\
LaBraM-Base~\cite{labram,eeg_bench} & Patch Transformer & 5.8M & 0.814 & 0.838 & -- & FT (ext.) \\
FEMBA-Base~\cite{femba} & Bidirec.\ Mamba & 47.7M & -- & 0.811 & 0.889 & FT (TUEG 21kh) \\
FEMBA-Huge~\cite{femba} & Bidirec.\ Mamba & 386M & -- & \textbf{0.818} & \textbf{0.901} & FT (TUEG 21kh) \\
\midrule
\multicolumn{7}{l}{\textit{Supervised baselines (TUAB labels only)}} \\
EEGNet~\cite{eegnet,kiessner2023} & Depthwise CNN & 9k & 0.796 & 0.824 & 0.913 & FT \\
Deep4Net~\cite{shallowconvnet} & Deep CNN & -- & 0.816 & -- & -- & FT \\
\midrule
\multicolumn{7}{l}{\textit{Published SSL models (WIN, TUAB-only pretraining)}} \\
ContraWR~\cite{contrawr} & CNN & 1.6M & 0.775 & -- & -- & Frozen \\
BIOT~\cite{biot} & FFT Transformer & 3.2M & 0.802 & -- & -- & Frozen \\
LaBraM-Base~\cite{labram} & Patch Transformer & 5.8M & \textbf{0.814} & -- & -- & Frozen \\
\midrule
\multicolumn{7}{l}{\textit{\cellcolor{ours}EB-JEPA (this work) --- frozen encoder, TUAB-only SSL}} \\
\cellcolor{ours}VICReg (fixed) & \cellcolor{ours}Conv1D & \cellcolor{ours}0.4M & \cellcolor{ours}-- & \cellcolor{ours}0.819 & \cellcolor{ours}0.900 & \cellcolor{ours}SSL Frozen \\
\cellcolor{ours}VICReg$^{*}$ + corruption & \cellcolor{ours}Conv1D & \cellcolor{ours}0.4M & \cellcolor{ours}0.770 & \cellcolor{ours}0.825 & \cellcolor{ours}0.904 & \cellcolor{ours}SSL Frozen \\
\cellcolor{ours}SIGReg + corruption & \cellcolor{ours}Conv1D & \cellcolor{ours}0.4M & \cellcolor{ours}0.775 & \cellcolor{ours}0.825 & \cellcolor{ours}0.904 & \cellcolor{ours}SSL Frozen \\
\cellcolor{ours}NaiveJEPA & \cellcolor{ours}Conv1D & \cellcolor{ours}0.4M & \cellcolor{ours}-- & \cellcolor{ours}0.779 & \cellcolor{ours}0.863 & \cellcolor{ours}SSL Frozen \\
\cellcolor{ours}VICReg$^{*}$ + spec 0.1$^{\dag}$ & \cellcolor{ours}Conv1D & \cellcolor{ours}0.4M & \cellcolor{ours}0.765 & \cellcolor{ours}\textbf{0.836} & \cellcolor{ours}0.887 & \cellcolor{ours}SSL Frozen \\
\midrule
\multicolumn{7}{l}{\textit{\cellcolor{ours}EB-JEPA (this work) --- fine-tuned}} \\
\cellcolor{ours}VICReg$^{*}$ + corruption $\to$ FT & \cellcolor{ours}Conv1D & \cellcolor{ours}0.4M & \cellcolor{ours}-- & \cellcolor{ours}0.812 & \cellcolor{ours}0.908 & \cellcolor{ours}Fine-tuned \\
\bottomrule
\multicolumn{7}{l}{$^{\dag}$Trained with $\lambda_{\rm inv}{=}1$ (default, outperforms $\lambda_{\rm inv}{=}25$ by $3$\,pp/$6\sigma$); 0.836 not reproducible due to eval-set coeff selection.}
\end{tabular}}
\end{table}

\paragraph{Discussion.}
\begin{itemize}[noitemsep]
  \item The 0.4M Conv1D encoder with SIGReg + corruption reaches 0.775 WIN BACC (= ContraWR,
    1.6M parameters) and 0.825 REC BACC within 20 training epochs without any labels.
  \item Per-recording, our fine-tuned probe reaches $0.812\pm0.004$ (3 seeds, final epoch),
    statistically matching supervised EEGNet ($0.824 \pm$ REC, FT).  SSL pretraining without
    labels thus \emph{ties} the supervised baseline at patient-level classification.
  \item The 0.836 REC BACC (spectral-0.1, $\lambda_{\rm inv}{=}1$) is not reproducible:
    it was eval-set-selected and single-seed.  The robust baseline is $0.819\pm0.004$
    (VICReg, $\lambda_{\rm inv}{=}1$, 3 seeds, corruption, no spectral).
    Notably, $\lambda_{\rm inv}{=}1$ \emph{outperforms} the paper-correct $\lambda_{\rm inv}{=}25$
    ($0.789\pm0.005$) by $3$\,pp --- a genuine EEG-specific finding, not a corrected bug.
  \item The 3.9\,pp gap to LaBraM-Base at WIN remains unexplained; tokeniser quality
    (VQ-codebook), model scale (5.8M vs.\ 0.4M), and pretraining duration (200 vs.\
    20 epochs) are the prime candidates.
\end{itemize}

%% ============================================================
\section{Conclusion and Future Work}
\label{sec:conclusion}
%% ============================================================

\subsection{Summary}

We presented EB-JEPA, a two-view energy-based SSL framework for clinical EEG with four
encoder families and multiple objective variants.
The primary contribution is \emph{diagnostic}: three systematic failure modes were
identified, quantified, and corrected during the hackathon.

\paragraph{Positive results.}
A 0.4M Conv1D encoder trained without labels for 20 epochs reaches 0.775 WIN BACC
(matching ContraWR) and 0.819 REC BACC (matching supervised EEGNet after fine-tuning).
NaiveJEPA, a minimal temporal-split predictive objective requiring no anti-collapse
regulariser, reaches 0.389 BACC on TUEV 6-class transfer \emph{with no event labels
during pretraining}, demonstrating cross-task generalisation.

\paragraph{Novel finding: EEG prefers weak VICReg invariance.}
A complete 4-point ablation ($\lambda_{\rm inv} \in \{1, 5, 10, 25\}$, endpoints
3-seed, intermediates single-seed) reveals a clear ordering:

\begin{center}\small
\begin{tabular}{cccc}
\toprule
$\lambda_{\rm inv}$ & BACC\textsubscript{REC} & AUROC & Seeds \\
\midrule
1 & $0.819 \pm .004$ & 0.900 & 3 \\
5 & $0.820$ & 0.893 & 1 \\
10 & $0.814$ & 0.900 & 1 \\
25 & $0.789 \pm .005$ & 0.882 & 3 \\
\bottomrule
\end{tabular}
\end{center}

The sweet spot is $\lambda \in \{1, 5\}$ (essentially tied at 0.819--0.820),
with a clear drop at $\lambda{=}10$ and a large drop at $\lambda{=}25$
($3\,\text{pp}$, ${\approx}6\sigma$ vs $\lambda{=}1$, confirmed 3-seed).
The ordering is not strictly monotone at the top ($\lambda{=}5 \approx \lambda{=}1$),
but the overall trend is unambiguous: for short-window (10\,s) clinical EEG
with a 0.4M conv encoder, anti-collapse terms (variance + covariance) should
dominate over view-invariance.  The paper-correct vision default ($\lambda{=}25$)
is the worst setting in the sweep.

\paragraph{Failure modes documented.}
\begin{enumerate}[noitemsep]
  \item \textbf{EEGPT channel-drop $\times$ channel-pool conflict}: the cross-attention
    pooling (19 channels $\to$ 1 token per time step) is structurally incompatible with
    channel-drop augmentation.  Removing channel-drop raises TUEV BACC from 0.182 to 0.358.
  \item \textbf{Spectral loss confound (D4 resolved)}:
    a full 4-point sweep ($w_{\rm spec} \in \{0.05, 0.10, 0.20, 0.30\}$) with
    $\lambda_{\rm inv}{=}1$ (our best setting) yields TUAB BACC $\in [0.795, 0.813]$
    vs.\ the no-spectral baseline of 0.819.
    Drops of 0.006--0.024\,pp are within single-seed noise.
    The earlier apparent gain (spectral-0.1 at 0.836) arose from a confounded evaluation:
    the coefficient was selected on the eval set and the run was single-seed.
    The accurate characterisation is \emph{neutral-to-mildly-harmful}; multi-seed
    confirmation would be needed to rule out even a small negative effect.
\end{enumerate}

\subsection{Open Scientific Questions}

\begin{description}[noitemsep]
  \item[D1] Is EEGPT collapse due to the channel-drop conflict, the VICReg bug, or both?
    Separate ablation needed: (i) no channel-drop + buggy VICReg, (ii) channel-drop +
    fixed VICReg, (iii) no channel-drop + fixed VICReg.
  \item[D2] What explains the 3.9\,pp gap to LaBraM at WIN?
    Candidates: tokeniser quality (VQ-codebook), model scale (5.8M vs.\ 0.4M),
    pretraining duration (200 vs.\ 20 epochs), data diversity.
  \item[D3] LogReg vs.\ MLP probe: the MLP inflates BACC by $\approx$0.01--0.02
    vs.\ LogReg; switching would enable direct comparison to SimCLR/MoCo literature.
  \item[D4] (\emph{Resolved}) Spectral loss: neutral-to-mildly-harmful.
    Full sweep ($w_{\rm spec} \in \{0.05, 0.10, 0.20, 0.30\}$, $\lambda_{\rm inv}{=}1$)
    shows drops of 0.006--0.024\,pp TUAB BACC, within single-seed noise.
    The earlier 0.836 (spec=0.1) was confounded: eval-set coefficient selection + single seed.
    No positive benefit exists at any tested coefficient.
  \item[D7] (\emph{Resolved}) Optimal $\lambda_{\rm inv}$ for EEG.
    Full sweep $\lambda \in \{1, 5, 10, 25\}$ complete.  Sweet spot: $\lambda \in \{1,5\}$
    (0.819--0.820 BACC), clear drop at $\lambda{=}10$ (0.814), large drop at
    $\lambda{=}25$ (0.789, $6\sigma$ below $\lambda{=}1$, 3-seed confirmed).
    Multi-seed runs at $\lambda \in \{5, 10\}$ would sharpen the boundary,
    but the conclusion is robust: EEG SSL requires substantially weaker
    view-invariance than the vision VICReg default.
  \item[D5] Aperiodic shortcut: does EB-JEPA's SIGReg representation exploit the 1/$f$
    spectral slope?  Apply~\cite{spectral_audit}'s phase-preserving Fourier intervention to
    frozen Conv1D + SIGReg to determine if BACC drops after flattening.
  \item[D6] Zero-shot abnormality scoring from NaiveJEPA energy: the MSE prediction error
    $E(\mathbf{x}) = \|\hat{\mathbf{z}}_{\rm tgt} - \mathbf{z}_{\rm tgt}\|^2$ computed per
    recording can serve as a direct abnormality score without any probe or labels.
    Computing AUROC from this energy alone is a strong test of world-model quality.
\end{description}

\subsection{Future Work}

\begin{enumerate}[noitemsep]
  \item \textbf{Neural tokeniser for LaBraM-style.}
    Training the VQ-codebook patch tokeniser~\cite{labram} (masked Fourier prediction on
    patches) and plugging it into our 1.2M Transformer backbone is the highest-priority
    architectural experiment.
  \item \textbf{Context ratio ablation for NaiveJEPA.}
    Try $r \in \{0.25, 0.50, 0.75\}$: smaller context = harder prediction = potentially
    richer representations (standard I-JEPA finding).
  \item \textbf{Longer pretraining.}
    All runs used 20 epochs.  LaBraM used $\approx$200 epochs on 2\,500h of diverse EEG.
    Even 40 epochs on TUAB-only is a cheap test of duration vs.\ capacity.
  \item \textbf{Mamba-based encoder.}
    FEMBA~\cite{femba} shows bidirectional Mamba scales linearly in sequence length vs.\
    $\mathcal{O}(N^2)$ for Transformers.  A small Mamba ($\approx$7.8M, FEMBA-Tiny)
    tested with VICReg is a natural next step.
  \item \textbf{NaiveJEPA + spectral on TUEV.}
    Stacking temporal masking (free anti-collapse) with spectral reg may combine
    complementary signals.
  \item \textbf{Multi-seed statistical significance.}
    All key findings rest on single-seed runs.  Rerunning the top 2--3 configs with 3 seeds
    would separate real signal from lucky initialisations.
\end{enumerate}

\subsection{Methodological Lessons}

\begin{enumerate}[noitemsep]
  \item \textbf{Always verify evaluation protocol.}
    WIN vs.\ REC-WIN is a 5\,pp systematic gap.  Comparing across protocols manufactures
    wins that vanish when the measurement is harmonised.
  \item \textbf{Loss coefficient deviations are invisible from convergence curves.}
    A $25{\times}$ underweighted invariance term still converges stably; the deviation
    is only detectable by auditing loss magnitudes --- and may turn out to be beneficial,
    as it did here ($\lambda{=}1$ beats $\lambda{=}25$ by 3\,pp on EEG).
  \item \textbf{Augmentation--architecture coupling.}
    Augmentation policies that are benign for Conv1D can be catastrophic for EEGPT
    (channel-drop + channel-pool).  Design or ablate augmentation per encoder family.
  \item \textbf{Best-epoch selection peeks at the test set.}
    Reporting the epoch with highest evaluation BACC uses the test set for model selection.
    Our fine-tuned probe appeared to reach 0.837 at its best epoch but only $0.812\pm0.004$
    at the final epoch (3 seeds).  The 2.5\,pp gap is pure peeking; it was exactly the margin
    that made SSL look like it beat EEGNet rather than tie it.
\end{enumerate}

\bigskip
\noindent\textbf{Code and data.}
All encoder architectures, training scripts, evaluation protocols, experiment logs,
and the science-check protocol are available in the project repository
(\texttt{Hack\_the\_worlds/}).

\bibliographystyle{plainnat}
\bibliography{references}

\end{document}
